%-*-tex-*-
% Copyright Michael J. Ferguson, INRS-Telecommunications
% All rights reserved. 

% ============ Tag Generation ==========
% These are the macros that generate the actual tags. They are used in 
% automatic equation numbering, section heads and any \auto...num{tag} form. 


% ============== Symbol that is dropped in when a tag is undefined ======
\def\undefsymbol{\hbox{\vrule width 1em height .7em}}

% ========== Tag types ============
% Two types of tags are now supported. They differ in that one form is
% immediate, and hence can refer only to things that have been defined while
% the other is delayed until the page is printed. These latter sometimes do
% migrate all the way out of internal boxes. 
% \ptag is the delayed form and \tag is the immediate form 

% page tags may turn out to be rather expensive. They are turned on through 
% the following form 
\newif\ifp@tag 
\def\pagetagson{\p@tagtrue} % this may require more work later 
\def\pagetagsoff{\p@tagfalse} % this is the default 


% ============ Tag macros =========
%    #1 is tag string 
%    #2 is tag value  (defined as the macro ":@#1")

\def\tag#1#2{\expandafter\xdef\csname :@#1\endcsname{#2}}

%   #1 is tag string 
%   #2 is tag value (defined as the macro ";@#1") 

\def\ptag#1#2{\expandafter\xdef\csname ;@#1\endcsname{#2}}

% \newttag makes an immediate \tag and writes it out immediately 

\def\newttag#1#2{\def\next{#1}\ifx \next\empty\tag{#1}{#2}\else
                  \tag{#1}{#2}\ift@agfilein \immediate\write\t@agfile 
               {\string\tag{#1}{#2}}\fi\proofmargin{#1 ==> #2}\fi}

% \newptag makes a page tag by writing it out to a file .. it is 
% not made immediately              


\def\newptag#1{\def\next{#1}\ifx \next\empty\relax\else
                  \ift@agfilein 
                   {\let\folio=\relax\xdef\next{\write\t@agfile 
          {\string\ptag{#1}{\folio}}}}\else\xdef\next{\relax}\fi
                   \fi\next}

\def\newtag#1#2{\newttag{#1}{#2}\ifp@tag\newptag{#1}\fi}


%undefined tag message

%\def\undeftagmessage#1{ <<<tag:#1 = undefined>>> } see english/french titles

% Tag reference macro
%    #1 is the string that was originally used in the tag
\def\r@ef#1#2{\def\doit{\ifundefined{#2#1}\writeterm{\undeftagmessage{#1}}{}{}\proofmargin
          {tag: #1 <-- undefined}\undefsymbol\else
           \proofmargin{ #1 <-- \csname #2#1\endcsname
             }\csname #2#1\endcsname
                       \fi}\relax
            \ifautoreference \doit \else#1\fi}
\def\q@ref#1#2{\ifautoreference \csname #2#1\endcsname\else#1\fi}

% ========= Normal Section Tag References ===========
\def\ref#1{\r@ef{#1}{:@}}
\def\quietref#1{\q@ref{#1}{:@}}

% ======== Page References ==========
\def\pageref#1{\r@ef{#1}{;@}}
\def\pagequietref#1{\q@ref{#1}{;@}}
\let\pref=\pageref
\let\qref=\quietref
\let\pqref=\pagequietref


% General autotag generation, used for equation numbers, figure numbers etc.
%    #1 = tag
%    #2 = counter
%    #3 = tag format
%    #4 = auto text
%    #5 = noauto text
%  does not generate text if \silenttrue 
\def\a@utotag#1#2#3#4#5{\ifautonumber
                           \global\advance #2 by 1 \relax 
                    \newtag{#1}{#3}\edef\t@ext{#4}\else \edef\t@ext{#5}\fi
                    \ifsilent\else\t@ext\fi}
